% ---------------------------------------------------------------------------
\section{Fundamentação Teórica}
\label{sec:fundamentacao}
% ---------------------------------------------------------------------------

Este capítulo apresenta os fundamentos teóricos que sustentam o estudo da geração automática de fases tridimensionais para jogos digitais. A exposição organiza-se de modo progressivo: parte-se da definição de fases, de \textit{Level Design} e da distinção conceitual entre organização lógica e organização espacial; em seguida, abordam-se os princípios da Geração Procedural de Conteúdo (PCG), sua vertente baseada em aprendizado de máquina (PCGML) e a caracterização dos Modelos de Linguagem de Grande Porte (LLMs), incluindo suas limitações de raciocínio e ancoragem espacial. Por fim, discutem-se a modelagem lógica de fases por meio da Teoria dos Grafos, a representação modular e geométrica de \textit{layouts} tridimensionais, as arquiteturas híbridas guiadas por restrições e os critérios formais de validade, navegabilidade e solvabilidade.

% ---------------------------------------------------------------------------
\subsection{Fases de Jogos, \textit{Level Design} e \textit{Layouts} 3D}
\label{sec:level_design}
% ---------------------------------------------------------------------------

\subsubsection{Fases, Níveis e \textit{Level Design}}
\label{sec:fases_niveis}

No contexto dos jogos digitais, uma fase --- também denominada nível (\textit{level}) ou mapa --- constitui a unidade espacial, temporal e interativa na qual as regras e mecânicas do sistema de regras do jogo se manifestam para o jogador \cite{tanagra, pcgbook}. Mais do que um cenário estático, uma fase opera como um contêiner estruturado para a experiência de jogo (\textit{gameplay}), delimitando o território de exploração, os recursos disponíveis, os obstáculos a serem transpostos e as condições de vitória ou transição \cite{tanagra, linden2014dungeons}. Em gêneros como ação e aventura, jogos de plataforma e \textit{Role-Playing Games} (RPGs), a qualidade da experiência interativa depende diretamente de como esses espaços são construídos e articulados \cite{dormans2010, linden2014dungeons}.

A disciplina responsável pelo planejamento, concepção e estruturação de fases é o \textit{Level Design} \cite{dormans2010, tanagra, totten2014architectural}. Seus objetivos fundamentais abrangem a condução da atenção e do deslocamento do jogador no espaço tridimensional, a introdução gradual de mecânicas, o balanceamento entre desafio e recompensa e a regulação do ritmo da experiência (\textit{pacing}) \cite{dormans2010, linden2013designing, linden2014dungeons, totten2014architectural}. Em fases do tipo masmorra (\textit{dungeon}), por exemplo, o \textit{Level Design} busca conciliar a liberdade de exploração espacial com um controle rigoroso sobre a progressão de desafios e quebra-cabeças (\textit{puzzles}), assegurando que o jogador não encontre obstáculos intransponíveis para o conjunto de habilidades ou recursos que detém em determinado instante da partida \cite{linden2014dungeons, shaker2016constructive}.

\subsubsection{\textit{Layout} de uma Fase}
\label{sec:layout_fase}

O \textit{layout} de uma fase refere-se ao arranjo estrutural e espacial dos elementos que compõem o ambiente jogável \cite{dormans2010, linden2014dungeons}. Em ambientes internos e masmorras, o \textit{layout} é classicamente decomposto em unidades arquitetônicas discretas, destacando-se as salas (ou câmaras), os corredores (ou vias de transição), os conectores (portas, portais, escadas e elevadores) e as regiões funcionais \cite{linden2014dungeons, shaker2016constructive, xu2025databasedrivenframework3dlevel}. As salas atuam tipicamente como arenas de interação concentrada --- abrigando combates, enigmas, recursos ou eventos narrativos ---, ao passo que os corredores estabelecem as ligações físicas entre essas áreas, modulando o tempo de travessia e a visibilidade entre setores adjacentes \cite{linden2013designing, shaker2016constructive}.

Além da geometria estrutural que delimita o espaço navegável (pisos, paredes, tetos e desníveis), o \textit{layout} compreende a disposição interna de elementos funcionais e estéticos, frequentemente denominados instalações ou objetos de cena (\textit{facilities} ou \textit{props}), bem como entidades dinâmicas, tais como personagens não jogáveis (NPCs), inimigos, armadilhas e itens coletáveis \cite{linden2014dungeons, xu2025constraint}. Em ambientes tridimensionais (3D), a especificação de um \textit{layout} envolve não apenas a distribuição planar em dois eixos, mas também a verticalidade, a sobreposição de múltiplos pavimentos, as relações de linha de visão (\textit{line of sight}) e as dimensões volumétricas exatas de cada componente no espaço contínuo \cite{xu2025constraint, xu2025databasedrivenframework3dlevel}.

\subsubsection{Organização Lógica e Organização Espacial}
\label{sec:organizacao_logica_espacial}

Um princípio teórico central na literatura de \textit{Level Design} consiste na distinção analítica entre duas estruturas complementares que coexistem em uma fase: a \textit{missão} (organização lógica) e o \textit{espaço} (organização espacial) \cite{dormans2010, linden2014dungeons}. A missão corresponde à sequência lógica ou parcialmente ordenada de tarefas, ações e objetivos que o jogador deve cumprir para concluir o nível, tais como derrotar um guardião, obter uma chave específica, desativar um bloqueio e alcançar a saída \cite{dormans2010, linden2013designing}. O espaço, por sua vez, refere-se à realização física e geométrica do ambiente virtual onde essas tarefas estão situadas, englobando as dimensões das salas, as distâncias métricas e a configuração arquitetônica das conexões \cite{dormans2010, linden2014dungeons}.

Embora o espaço físico acomode a missão e esta seja mapeada sobre aquele, as duas estruturas são conceitualmente independentes e regidas por princípios distintos \cite{dormans2010}. Uma missão estritamente linear --- na qual os objetivos precisam ser cumpridos em uma ordem única e fixa --- pode ser distribuída sobre um espaço físico não linear, como um arranjo em estrela (\textit{hub-and-spoke}) ou em anel, exigindo que o jogador retorne a áreas previamente visitadas à medida que desbloqueia novos acessos \cite{dormans2010}. Inversamente, uma missão não linear, dotada de múltiplos desafios paralelos ou caminhos alternativos, pode ser acomodada em uma geometria espacial simples \cite{dormans2010, linden2013designing}. Essa dissociação teórica evidencia que a coerência da progressão de um nível reside primariamente em sua estrutura lógica de dependências, enquanto sua viabilidade física e navegacional depende da correta realização geométrica no espaço tridimensional \cite{dormans2010, linden2013designing, linden2014dungeons}.

% ---------------------------------------------------------------------------
\subsection{Geração Procedural de Conteúdo}
\label{sec:pcg}
% ---------------------------------------------------------------------------

\subsubsection{Conceito e Objetivos da PCG}
\label{sec:conceito_pcg}

A Geração Procedural de Conteúdo (PCG --- \textit{Procedural Content Generation}) é definida na literatura como a criação algorítmica de conteúdos para jogos digitais de maneira autônoma ou com intervenção humana limitada e indireta \cite{pcgbook, summerville2018procedural, maleki2024procedural}. O escopo do termo ``conteúdo'' abrange praticamente todos os artefatos que integram um jogo, incluindo níveis, mapas, terrenos, vegetação, texturas, efeitos sonoros, narrativas, missões (\textit{quests}) e conjuntos de regras, distinguindo-se do comportamento em tempo de execução dos agentes de inteligência artificial que controlam NPCs \cite{pcgbook, summerville2018procedural, maleki2024procedural}.

Historicamente introduzida em títulos pioneiros da década de 1980, como \textit{Rogue} (1980) e \textit{Elite} (1984), para contornar severas restrições de memória e armazenamento por meio da síntese algorítmica em tempo de execução \cite{pcgbook, summerville2018procedural, araujo2018geracao}, a PCG consolidou-se tanto na indústria quanto na pesquisa acadêmica em função de múltiplos objetivos \cite{linden2014dungeons, pcgbook}. Dentre as principais motivações para o seu emprego destacam-se: (i) a automação e a redução do tempo e do custo financeiro associados à autoria manual de ativos em larga escala; (ii) o aumento da variabilidade e da diversidade dos ambientes gerados, ampliando o fator de rejogabilidade (\textit{replayability}); (iii) o suporte à adaptabilidade centrada no jogador, permitindo ajustar dinamicamente os desafios às habilidades observadas; e (iv) a exploração de novas estéticas e mecânicas emergentes \cite{linden2014dungeons, pcgbook, summerville2018procedural}.

A avaliação e o projeto de geradores procedurais apoiam-se em propriedades técnicas fundamentais, notadamente a reprodutibilidade, a expressividade e a controlabilidade \cite{tanagra, smith2010analyzing, linden2014dungeons, maleki2024procedural}. A reprodutibilidade decorre frequentemente do uso de Geradores de Números Pseudoaleatórios (PRNGs), nos quais a especificação de um valor inicial --- denominado semente (\textit{seed}) --- permite reconstruir de maneira determinística a mesma sequência de decisões algorítmicas \cite{maleki2024procedural}. A expressividade, caracterizada pela análise do espaço expressivo (\textit{expressive range}), quantifica e mapeia a variedade, a distribuição e as propriedades estilísticas ou estruturais dos artefatos que o gerador é capaz de produzir no espaço de design \cite{smith2010analyzing, linden2013designing}. A controlabilidade, por sua vez, reflete a capacidade do projetista ou do sistema de direcionar intencionalmente as características do conteúdo produzido por meio de parâmetros, regras ou restrições \cite{linden2014dungeons, shaker2016constructive}. Um desafio central no projeto de geradores consiste em equilibrar a tensão entre controle e diversidade: restrições excessivamente rígidas podem limitar a variabilidade dos resultados, ao passo que a ausência de controle eleva o risco de geração de artefatos inválidos ou incoerentes \cite{smith2010analyzing, linden2014dungeons}.

\subsubsection{Principais Abordagens de PCG}
\label{sec:abordagens_pcg}

Uma classificação recorrente na literatura de PCG organiza os métodos generativos em três grandes paradigmas operacionais: métodos construtivos, métodos baseados em busca (\textit{search-based}) e métodos baseados em restrições ou solucionadores (\textit{solver-based} / \textit{constraint-based}) \cite{togelius2011search, pcgbook, summerville2018procedural, maleki2024procedural}.

Os métodos construtivos caracterizam-se por gerar o conteúdo em uma única passagem sequencial de passos algorítmicos, com tempo de execução previsível e sem realizar ciclos iterativos de avaliação e regeneração sobre o artefato produzido \cite{shaker2016constructive, summerville2018procedural}. Exemplos clássicos incluem funções de ruído coerente (como o ruído de Perlin para síntese de terrenos e texturas), autômatos celulares, algoritmos de particionamento espacial e gramáticas generativas \cite{linden2014dungeons, shaker2016constructive, maleki2024procedural}. Importa ressaltar que procedimentos construtivos podem fornecer garantias estruturais formais sobre o conteúdo gerado (como ausência de sobreposição ou conectividade garantida) quando suas regras forem rigorosamente formuladas para preservar tais invariantes em cada etapa da construção \cite{shaker2016constructive}. Por outro lado, caso dependam de decisões puramente estocásticas desprovidas de verificação antecipada, o determinismo algorítmico assegura apenas a repetibilidade da saída a partir de uma semente fixa, não prevenindo falhas estruturais ou becos sem saída \cite{linden2014dungeons, shaker2016constructive}.

Os métodos baseados em busca (\textit{search-based PCG}) enquadram-se como um caso particular do paradigma gerar-e-testar (\textit{generate-and-test}) guiado por otimização \cite{togelius2011search, pcgbook, maleki2024procedural}. Em vez de aceitar imediatamente o artefato sintetizado ou simplesmente descartá-lo de forma binária, algoritmos baseados em busca atribuem a cada instância candidata um valor numérico contínuo por meio de uma função de avaliação ou aptidão (\textit{fitness function}), utilizando esse escore para guiar iterativamente a modificação e a seleção de novas soluções no espaço de busca \cite{togelius2011search, linden2014dungeons, araujo2018geracao, maleki2024procedural}. Essa categoria emprega amplamente técnicas de computação evolutiva (como Algoritmos Genéticos) \cite{togelius2011search, linden2014dungeons, maleki2024procedural}, além de meta-heurísticas como têmpera simulada (\textit{Simulated Annealing}), otimização por enxame de partículas (PSO) e Busca em Árvore de Monte Carlo (MCTS) \cite{maleki2024procedural, xu2025constraint}.

Os métodos baseados em restrições (\textit{constraint-based} ou \textit{solver-based PCG}), por sua vez, especificam declarativamente as propriedades e relações obrigatórias que o artefato final deve satisfazer, delegando a construção a solucionadores matemáticos ou lógicos (\textit{constraint solvers}) \cite{tanagra, smith2011answer, linden2014dungeons, pcgbook, summerville2018procedural}. Destacam-se nesse paradigma a Programação por Conjuntos de Respostas (\textit{Answer Set Programming} --- ASP), que formaliza regras e restrições de design como predicados lógicos satisfatíveis \cite{smith2011answer}, solucionadores de satisfatibilidade módulo teorias (SMT) e solucionadores de Problemas de Satisfação de Restrições (CSP) \cite{tanagra, chen2025narrative}. Integram também essa vertente os algoritmos baseados em propagação local de restrições de adjacência inspirados em mecânica estatística, com destaque para o \textit{WaveFunctionCollapse} (WFC) \cite{karth2017wavefunctioncollapse, maleki2024procedural}. Nessa abordagem, o foco desloca-se da definição de passos procedurais imperativos para a formulação matemática e lógica do espaço de design admissível \cite{tanagra, smith2011answer, linden2014dungeons}.

\subsubsection{Geração Procedural de Fases}
\label{sec:pcg_fases}

A aplicação de PCG à criação de fases e níveis constitui um dos tópicos mais investigados na área, apresentando exigências técnicas particulares que a diferenciam da geração de ativos puramente cosméticos \cite{summerville2018procedural, maleki2024procedural}. A literatura distingue conteúdos cosméticos (como texturas e efeitos visuais, cuja alteração não modifica o resultado das ações do jogador) de conteúdos funcionais, que afetam diretamente a mecânica e a dinâmica da partida \cite{summerville2018procedural}. Fases de jogos são artefatos intrinsecamente funcionais: uma fase visualmente atraente, mas que apresente falhas topológicas ou bloqueios estruturais que impeçam sua conclusão, é considerada inválida para o propósito do jogo \cite{summerville2018procedural, linden2014dungeons}.

Na geração procedural de níveis e masmorras, um método típico envolve três componentes conceituais: (1) um modelo representacional abstrato que descreve a estrutura global da fase; (2) um método para construir esse modelo representacional; e (3) um método para sintetizar a geometria efetiva da fase a partir do modelo abstrato \cite{linden2014dungeons, shaker2016constructive}. Diversas famílias algorítmicas foram desenvolvidas para essas etapas. Algoritmos de particionamento de espaço, como o Particionamento Binário do Espaço (\textit{Binary Space Partitioning} --- BSP), \textit{quadtrees} (em 2D) e \textit{octrees} (em 3D), subdividem recursivamente o domínio espacial em células disjuntas, garantindo estruturalmente que salas alocadas em partições distintas não se sobreponham \cite{shaker2016constructive}. Em contrapartida, abordagens baseadas em agentes escavadores (\textit{agent-based dungeon growing}) operam em escala local, deslocando um agente virtual que escava túneis e posiciona salas de maneira passo a passo; embora produzam configurações orgânicas, agentes sem mecanismos de inspeção antecipada (\textit{look-ahead}) estão sujeitos a gerar salas sobrepostas e corredores sem saída \cite{shaker2016constructive}.

De modo complementar, autômatos celulares operam sobre grades discretas aplicando regras de transição baseadas na vizinhança local (como as vizinhanças de Moore ou de von Neumann) para formar cavernas naturais, mas não asseguram, por si sós, a conectividade global entre câmaras isoladas, exigindo etapas adicionais de verificação e abertura de passagens \cite{linden2014dungeons, shaker2016constructive}. Gramáticas generativas --- incluindo gramáticas de grafos e gramáticas de formas (\textit{shape grammars}) --- expandem um axioma inicial por meio de regras de produção e reescrita, permitindo capturar princípios de progressão e hierarquia em níveis de ação e aventura \cite{dormans2010, linden2013designing, shaker2016constructive}. Apesar da maturidade dessas técnicas em domínios bidimensionais, a transição para a geração de fases tridimensionais complexas permanece um desafio aberto na literatura devido ao aumento expressivo das restrições geométricas, de navegação vertical e de acoplamento espacial \cite{linden2014dungeons, maleki2024procedural, xu2025constraint}.

% ---------------------------------------------------------------------------
\subsection{PCG Baseada em Aprendizado de Máquina e Uso de LLMs}
\label{sec:pcgml_llms}
% ---------------------------------------------------------------------------

\subsubsection{\textit{Procedural Content Generation via Machine Learning} --- PCGML}
\label{sec:pcgml}

A Geração Procedural de Conteúdo via Aprendizado de Máquina (\textit{Procedural Content Generation via Machine Learning} --- PCGML) é definida formalmente como a geração de conteúdos de jogos diretamente por modelos estatísticos ou de aprendizado de máquina treinados sobre conjuntos de conteúdos preexistentes de jogos \cite{summerville2018procedural}. A distinção fundamental entre o PCGML e as abordagens clássicas baseadas em busca ou em solucionadores reside no mecanismo gerador: enquanto a busca evolucionária pode eventualmente empregar um modelo neural apenas como função de avaliação para pontuar candidatos sintetizados por outros meios, no PCGML o modelo treinado internaliza a distribuição estatística do \textit{corpus} de entrada e amostra diretamente novas instâncias a partir de sua inferência \cite{summerville2018procedural}.

A taxonomia clássica de PCGML organiza os métodos segundo a representação subjacente dos dados (sequências unidimensionais, grades bidimensionais/tridimensionais e grafos) e o paradigma de modelagem empregado \cite{summerville2018procedural}. Entre as técnicas consolidadas figuram: (i) modelos probabilísticos de contagem de frequência, como modelos de $n$-gramas, Cadeias de Markov Multidimensionais (MdMCs) e Campos Aleatórios de Markov (MRFs); (ii) métodos de redução de dimensionalidade e fatoração de matrizes (como PCA); (iii) Redes Neurais Recorrentes, especialmente redes \textit{Long Short-Term Memory} (LSTMs), utilizadas na geração sequencial de fases baseadas em blocos; (iv) \textit{Autoencoders} e \textit{Variational Autoencoders} (VAEs), que mapeiam fases em espaços latentes contínuos para fins de interpolação e reparo de segmentos; e (v) Redes Adversárias Generativas (GANs) \cite{summerville2018procedural, maleki2024procedural}.

No âmbito do aprendizado por reforço aplicado à geração (PCGRL), o paradigma difere do treinamento supervisionado sobre \textit{corpora} estáticos \cite{khalifa2020pcgrl}: o agente gerador interage dinamicamente com um ambiente de desenvolvimento de fases, modificando células ou inserindo peças por meio de ações discretas e sendo otimizado para maximizar uma função de recompensa cumulativa que quantifica propriedades de jogabilidade e solvabilidade \cite{khalifa2020pcgrl, maleki2024procedural}.

É conceitualmente relevante distinguir a utilização de Modelos de Linguagem de Grande Porte (LLMs) em relação à definição estrita de PCGML formulada por Summerville et al. \cite{summerville2018procedural}. Quando um LLM fundacional é empregado por meio de engenharia de \textit{prompt} em regime de zero-shot ou few-shot, o modelo não foi treinado sobre um \textit{corpus} de fases de jogos específico, mas opera por meio de aprendizado em contexto (\textit{in-context learning}) suportado pelo conhecimento estrutural e semântico adquirido durante o pré-treinamento geral em linguagem natural e código \cite{brown2020language, Gallotta_2024}. Por outro lado, caso um modelo de linguagem seja expressamente submetido a ajuste fino (\textit{fine-tuning}) sobre conjuntos de dados de fases (como representações em texto, ASCII ou grafos de níveis), sua operação alinha-se estritamente à definição clássica de PCGML \cite{summerville2018procedural, todd2023level}.

Embora o aprendizado a partir de exemplos atenue a necessidade de codificação manual de regras estilísticas, a literatura aponta duas limitações estruturais no paradigma de PCG baseado em modelos conexionistas \cite{summerville2018procedural, maleki2024procedural, her2025zeroshot3dmapgeneration}: a escassez de dados (\textit{data scarcity}), decorrente do número restrito de fases digitalizadas e da ausência de representações universais entre diferentes jogos \cite{summerville2018procedural}; e a dificuldade de garantir propriedades funcionais globais, visto que modelos puramente estatísticos tendem a reproduzir padrões visuais plausíveis localmente, mas não oferecem garantias determinísticas de conectividade, acessibilidade ou solvabilidade global sem módulos externos de validação \cite{summerville2018procedural, maleki2024procedural}.

\subsubsection{Modelos de Linguagem de Grande Porte Aplicados à Geração de Conteúdo}
\label{sec:llms_pcg}

Os Modelos de Linguagem de Grande Porte (\textit{Large Language Models} --- LLMs) constituem uma classe de modelos fundacionais de aprendizado profundo fundamentados na arquitetura \textit{Transformer} \cite{vaswani2017attention, Gallotta_2024}. Historicamente, os modelos neurais sequenciais precursores em Processamento de Linguagem Natural (PLN), tais como Redes Neurais Recorrentes (RNNs) e LSTMs, enfrentavam severos entraves para capturar dependências de longo alcance devido ao fenômeno do esvanecimento ou explosão de gradientes (\textit{vanishing/exploding gradients}) durante o treinamento por retropropagação ao longo do tempo \cite{bengio1994learning}. A arquitetura \textit{Transformer} superou essa barreira teórica ao prescindir da recorrência sequencial e introduzir o mecanismo de autoatenção (\textit{self-attention}) \cite{vaswani2017attention}. O mecanismo de autoatenção calcula diretamente coeficientes de ponderação entre todos os pares de elementos (\textit{tokens}) da sequência de entrada, viabilizando o processamento altamente paralelizado e permitindo correlacionar informações estruturais e semânticas independentemente de sua distância posicional no texto \cite{vaswani2017attention, Gallotta_2024}.

Na vertente generativa preponderante, os modelos autorregressivos causais operam estimando a distribuição de probabilidade condicional do próximo \textit{token} $x_t$ com base no histórico de \textit{tokens} precedentes:
\begin{equation}
    P(x_1, x_2, \dots, x_T) = \prod_{t=1}^{T} P(x_t \mid x_1, x_2, \dots, x_{t-1})
\end{equation}
O ciclo de desenvolvimento desses modelos envolve tipicamente o pré-treinamento autossupervisionado em larga escala sobre extensos \textit{corpora} de textos e linguagens de programação, seguido por ajuste fino supervisionado para seguimento de instruções (\textit{instruction tuning}) e alinhamento por preferências (como RLHF) \cite{brown2020language, Gallotta_2024, her2025zeroshot3dmapgeneration}. Em modelos com bilhões de parâmetros, destaca-se a capacidade de aprendizado em contexto (\textit{In-Context Learning} --- ICL), que possibilita a resolução de novas tarefas a partir exclusivamente de especificações no \textit{prompt} em regime \textit{zero-shot} ou condicionada por poucos exemplos demonstrativos (\textit{few-shot}), sem modificação dos pesos sinápticos da rede \cite{brown2020language}. Estruturações avançadas de \textit{prompting}, como Cadeia de Pensamento (\textit{Chain-of-Thought} --- CoT) e a incorporação de Geração Aumentada por Recuperação (\textit{Retrieval-Augmented Generation} --- RAG), estendem essa capacidade permitindo decomposição analítica de etapas e consulta a fontes de dados estruturados \cite{Gallotta_2024, tang2025unrealllm, wang2026designerlyloop}.

Na literatura de PCG, o emprego de LLMs diversificou-se em múltiplas funções especializadas \cite{Gallotta_2024, maleki2024procedural}:
\begin{enumerate}
    \item \textbf{Gerador direto de representações e conteúdo:} síntese de artefatos textuais como diálogos e narrativas, bem como mapas bidimensionais codificados diretamente em texto ou caracteres ASCII (como grades para \textit{Sokoban} ou \textit{Super Mario Bros.}) \cite{todd2023level, nasir2024word2worldgeneratingstoriesworlds};
    \item \textbf{Planejador lógico e semântico de alto nível:} geração de estruturas conceituais abstratas, incluindo árvores de missões, grafos de progressão lógica e atribuição contextual de temas narrativos a salas e regiões \cite{Gallotta_2024, xu2025databasedrivenframework3dlevel, wang2026designerlyloop};
    \item \textbf{Sintetizador de parâmetros e código para ferramentas procedurais:} produção de rotinas em linguagens de programação (como scripts para motores ou modeladores 3D), parametrização de geradores construtivos e formulação de restrições para solucionadores matemáticos \cite{hu2024scenecraft, tang2025unrealllm, xu2025constraint, yang2026codeasroomgenerating3drooms};
    \item \textbf{Assistente de design em sistemas de iniciativa mista:} atuação colaborativa na interpretação de comandos do designer, oferecendo sugestões estruturais e refinando especificações conceituais \cite{Gallotta_2024, bandeali2026generative, wang2026designerlyloop}.
\end{enumerate}

\subsubsection{Geração de Representações Estruturadas}
\label{sec:representacoes_estruturadas}

Quando modelos de linguagem são integrados a sistemas computacionais externos --- tais como motores gráficos, otimizadores matemáticos ou verificadores formais ---, a emissão de respostas em linguagem natural livre torna-se inadequada devido à sua ambiguidade intrínseca e à suscetibilidade a falhas sintáticas na análise automatizada (\textit{parsing}) \cite{xu2025constraint, regmi2025system}. Para garantir a interoperabilidade entre a inferência probabilística do LLM e os procedimentos algorítmicos determinísticos, a literatura preconiza a geração de representações estruturadas intermediárias \cite{hu2024scenecraft, tang2025unrealllm, xu2025constraint, chen2025narrative, wang2026designerlyloop}.

Sob o ponto de vista conceitual e de engenharia de software, é indispensável diferenciar o \textit{modelo conceitual} do \textit{formato de serialização} \cite{xu2025constraint, xu2025databasedrivenframework3dlevel, chen2025narrative}:
\begin{itemize}
    \item O \textbf{modelo conceitual} (ou estrutura abstrata de informação) é a formalização lógica ou matemática que organiza o domínio do problema. Exemplos compreendem grafos de cena e de progressão \cite{xu2025databasedrivenframework3dlevel, wang2026designerlyloop}, conjuntos de triplas relacionais do tipo \texttt{[Objeto-Relação-Objeto]} \cite{chen2025narrative}, hierarquias de componentes de cena, linguagens de descrição de domínio específico (DSL) como a GLDL (\textit{Game Level Description Language}) \cite{xu2025constraint}, ou sínteses procedurais via código executável \cite{hu2024scenecraft, yang2026codeasroomgenerating3drooms};
    \item O \textbf{formato de serialização} é o padrão textual ou binário adotado para codificar, armazenar e transmitir esse modelo conceitual entre diferentes subsistemas \cite{sun2024roomplanner, xu2025databasedrivenframework3dlevel}. O formato JSON (\textit{JavaScript Object Notation}), bem como esquemas em XML ou YAML, constitui a representação hegemônica por estruturar dados hierárquicos em pares chave-valor e arranjos vetoriais com suporte a validação rigorosa de sintaxe por meio de esquemas formais (\textit{JSON Schema}) \cite{tang2025unrealllm, sun2024roomplanner, xu2025databasedrivenframework3dlevel}. Abordagens pontuais na literatura, a exemplo de Öcal et al. \cite{ocal2024sceneteller}, adaptaram experimentalmente a sintaxe de folhas de estilo (CSS) para declarar parâmetros de caixas delimitadoras; não obstante, o padrão JSON consolida-se como a convenção primária de serialização para intercâmbio de dados estruturados.
\end{itemize}

O formato de serialização não se confunde, portanto, com a lógica do nível: ele é o meio físico-textual de transmissão que transporta a abstração formal da fase para que ferramentas computacionais possam interpretá-la e validá-la de modo determinístico \cite{tang2025unrealllm, xu2025databasedrivenframework3dlevel}.

% ---------------------------------------------------------------------------
\subsection{Raciocínio e Ancoragem Espacial em LLMs}
\label{sec:spatial_grounding}
% ---------------------------------------------------------------------------

\subsubsection{Raciocínio Espacial}
\label{sec:raciocinio_espacial}

O raciocínio espacial compreende a capacidade de representar formalmente entidades em um domínio multidimensional, inferir suas propriedades geométricas e deduzir os desdobramentos de suas transformações e arranjos mútuos \cite{cohn2008qualitative, sun2024roomplanner, xu2025constraint}. No domínio do \textit{Level Design} e da síntese de ambientes virtuais 3D, o raciocínio espacial subdivide-se entre relações qualitativas (topológicas e mereológicas) \cite{cohn2008qualitative} e relações quantitativas métricas contínuas \cite{xu2025constraint, chen2025narrative}.

Na formulação da organização espacial de uma fase, identificam-se propriedades e relações geométricas essenciais \cite{cohn2008qualitative, sun2024roomplanner, xu2025constraint, chen2025narrative}:
\begin{itemize}
    \item \textbf{Posição:} localização absoluta ou relativa de uma entidade no espaço $\mathbb{R}^3$, expressa em relação à origem global ou a um ponto de referência local;
    \item \textbf{Distância e proximidade:} separação métrica euclidiana entre entidades, determinando adjacência funcional, raio de alcance ou zonas de segurança;
    \item \textbf{Direção versus Orientação:} sob o rigor do raciocínio espacial \cite{cohn2008qualitative}, a \textit{direção} constitui uma relação espacial relativa entre duas entidades (indicando, por exemplo, que a sala $B$ está situada ao norte, a leste ou acima da sala $A$), enquanto a \textit{orientação} refere-se ao alinhamento angular intrínseco de uma entidade (sua \textit{pose} ou vetor de mira frontal --- \textit{facing} --- no espaço, determinado por suas rotações em torno dos eixos locais);
    \item \textbf{Adjacência e conectividade:} contato geométrico ou comunicação de passagem entre regiões contíguas, tais como salas que compartilham portais comuns;
    \item \textbf{Contenção:} exigência topológica e volumétrica de que determinado conjunto de elementos (como itens ou mobílias) esteja integralmente situado dentro do perímetro delimitador de uma sala ou compartimento hospedeiro;
    \item \textbf{Interseção e ocupação volumétrica:} verificação de compartilhamento de pontos no espaço tridimensional entre dois ou mais corpos físicos.
\end{itemize}

A preservação concorrente da consistência de todas essas relações ao longo da alocação de dezenas de elementos arquitetônicos e funcionais constitui um problema combinatório e contínuo de alta complexidade \cite{xu2025constraint, feng2026repurposing3dgenerativemodel}.

\subsubsection{\textit{Spatial Grounding}}
\label{sec:conceito_spatial_grounding}

O problema de ancoragem semântica (\textit{symbol grounding problem}), formulado originalmente por Harnad \cite{harnad1990symbol}, investiga como símbolos abstratos manipulados puramente com base em regras sintáticas podem adquirir significado intrínseco e correspondência funcional com entidades concretas de um mundo físico ou simulado. No âmbito dos sistemas generativos aplicados a ambientes 3D, a \textit{ancoragem espacial} (\textit{spatial grounding}) designa a capacidade de mapear comandos linguísticos e especificações conceituais (como ``posicionar uma mesa no centro da câmara'' ou ``vincular o corredor de acesso à ala norte'') em coordenadas métricas coerentes, dimensões tridimensionais compatíveis e restrições espaciais verificáveis \cite{harnad1990symbol, hu2024scenecraft, sun2024roomplanner, chen2025narrative}.

Um componente crítico da ancoragem espacial consiste na interpretação correta de \textit{affordances}. Na definição clássica da psicologia ecológica de Gibson \cite{gibson1979ecological}, uma \textit{affordance} traduz as possibilidades reais de ação que os elementos do ambiente físico oferecem a um agente (por exemplo, uma porta aberta oferece passagem, uma superfície horizontal firme oferece suporte para caminhar ou sentar, e uma parede oferece bloqueio ou proteção) \cite{gibson1979ecological}. A ancoragem espacial requer que o sistema não apenas reconheça a coocorrência semântica entre termos linguísticos, mas garanta que as configurações geométricas resultantes respeitem as propriedades funcionais de uso: corredores devem manter altura e largura livres para circulação de agentes, portas devem possuir desobstrução frontal para acesso e objetos de apoio devem respeitar o contato estável de superfícies no espaço tridimensional \cite{sun2024roomplanner, xu2025constraint}.

\subsubsection{Limitações dos LLMs em Tarefas Espaciais Tridimensionais}
\label{sec:limitacoes_espaciais_llms}

No contexto dos modelos generativos de linguagem, define-se por alucinação a síntese de respostas que apresentam correção superficial ou gramatical, mas que violam a coerência factual, contêm premissas contraditórias ou desrespeitam restrições explícitas do domínio \cite{Gallotta_2024, xu2025constraint, bandeali2026generative}. Quando os LLMs são incumbidos de determinar diretamente atributos geométricos contínuos em ambientes 3D --- tais como coordenadas cartesianas $(x, y, z)$, caixas delimitadoras e ângulos de rotação para múltiplos elementos interconectados ---, esse comportamento manifesta-se sob a forma de \textit{alucinação espacial} ou inconsistência métrica \cite{Gallotta_2024, sun2024roomplanner, yang2026codeasroomgenerating3drooms, feng2026repurposing3dgenerativemodel}.

A fundamentação teórica dessa limitação decorre de características constitutivas dos modelos de linguagem \cite{Gallotta_2024, hu2024scenecraft, sun2024roomplanner, feng2026repurposing3dgenerativemodel}:
\begin{enumerate}
    \item \textbf{Ausência de modelagem física e geométrica na inferência:} LLMs são treinados com base na verossimilhança estatística de sequências de \textit{tokens} textuais, sem acesso nativo a um motor de cálculo vetorial ou representação interna de espaço euclidiano tridimensional durante o cálculo de autoatenção \cite{Gallotta_2024, sun2024roomplanner};
    \item \textbf{Severa compressão de informação geométrica:} ao representar geometrias poligonais complexas exclusivamente por rótulos simbólicos e tuplas numéricas discretizadas, omitem-se as interações de superfície contínuas e as relações de contato entre sólidos \cite{feng2026repurposing3dgenerativemodel};
    \item \textbf{Deriva autorregressiva e acúmulo de erros métricos:} em processos de geração sequencial ou em planejamentos em passo único (\textit{single-step}), pequenos desvios numéricos cometidos nas primeiras coordenadas preditas propagam-se recursivamente pelo contexto, degradando a consistência geométrica global conforme o número de elementos cresce no ambiente \cite{sun2024roomplanner, feng2026repurposing3dgenerativemodel, yang2026codeasroomgenerating3drooms}.
\end{enumerate}

Diversos estudos empíricos recentes ilustram concretamente essas falhas:
\begin{itemize}
    \item No sistema \textit{SceneCraft}, Hu et al. \cite{hu2024scenecraft} evidenciaram que, ao gerar scripts em Python para modeladores 3D como o Blender, os LLMs conseguem estruturar a hierarquia de objetos, mas cometem graves imprecisões nas coordenadas numéricas, gerando colisões e desalinhamentos severos que exigem otimizadores numéricos acoplados para corrigir a cena;
    \item No framework \textit{RoomPlanner}, Sun et al. \cite{sun2024roomplanner} constataram empiricamente sobre o conjunto de dados Matterport3D que a solicitação direta de caixas delimitadoras a LLMs resulta em taxas substanciais de interpenetração de móveis e posicionamento de elementos fora das paredes dos cômodos;
    \item No estudo \textit{Code-as-Room}, Yang et al. \cite{yang2026codeasroomgenerating3drooms} reportaram que a síntese de recintos 3D via código por LLMs padece de alucinações espaciais e violações funcionais de circulação quando desprovida de verificação externa de navegabilidade.
\end{itemize}

As falhas observadas agrupam-se em quatro categorias recorrentes \cite{sun2024roomplanner, xu2025constraint, feng2026repurposing3dgenerativemodel, yang2026codeasroomgenerating3drooms}: (i) \textbf{colisões e sobreposições volumétricas (\textit{overlapping} / \textit{clipping})}, nas quais malhas sólidas ocupam o mesmo volume simultaneamente; (ii) \textbf{violações de limites e suporte (\textit{out-of-bounds} e objetos flutuantes)}, em que elementos atravessam paredes externas ou ficam suspensos sem contato com o piso; (iii) \textbf{erros de rotação e alinhamento}, como portas e armários posicionados de frente para paredes cegas; e (iv) \textbf{obstruções de circulação}, nas quais objetos bloqueiam corredores vitais e fragmentam o espaço navegável \cite{sun2024roomplanner, xu2025constraint, xu2025databasedrivenframework3dlevel}.

Tais constatações consolidam a compreensão de que, embora os LLMs demonstrem notável competência para o raciocínio semântico e planejamento conceitual, a determinação direta de layouts tridimensionais estritamente por inferência de linguagem é inerentemente falha, demandando arquiteturas estruturadas que deleguem a instanciação geométrica e a validação espacial a componentes algorítmicos dedicados \cite{hu2024scenecraft, sun2024roomplanner, xu2025constraint, xu2025databasedrivenframework3dlevel, yang2026codeasroomgenerating3drooms}.

% ---------------------------------------------------------------------------
\subsection{Representação Lógica de Fases por Grafos}
\label{sec:grafos}
% ---------------------------------------------------------------------------

\subsubsection{Conceitos Fundamentais de Teoria dos Grafos}
\label{sec:conceitos_grafos}

A Teoria dos Grafos fornece a base matemática formal para modelar relações discretas e conectividades estruturais em sistemas complexos \cite{cormen2009introduction, summerville2018procedural}. Formalmente, um grafo é definido como um par ordenado $G = (V, E)$, no qual $V$ consiste em um conjunto finito e não vazio de vértices (ou nós) e $E$ representa um conjunto de arestas (ou arcos) que conectam pares de vértices pertencentes a $V$ \cite{cormen2009introduction, maleki2024procedural}.

Quanto à orientação de suas arestas, um grafo pode ser não direcionado --- em que as arestas indicam conexões simétricas bidirecionais entre pares de vértices $\{u, v\}$ ---, ou direcionado (dígrafo), quando cada aresta consiste em um par ordenado $(u, v) \in E$ orientado do vértice de origem $u$ para o vértice de destino $v$ \cite{cormen2009introduction, dormans2010}. Em dígrafos, distingue-se o grau de entrada ($d^-(v)$, número de arestas que incidem em $v$) do grau de saída ($d^+(v)$, número de arestas que partem de $v$) \cite{cormen2009introduction}. Um grafo ponderado associa valores numéricos reais às suas arestas (como distâncias métricas, tempos de travessia ou dificuldades), enquanto um grafo rotulado ou de atributos vincula tipos, símbolos e dicionários de propriedades tanto a vértices quanto a arestas \cite{dormans2010, linden2013designing, linden2014dungeons, chen2025narrative}.

Sobre essas estruturas, fundamentam-se propriedades topológicas cruciais \cite{cormen2009introduction, linden2014dungeons, summerville2018procedural}:
\begin{itemize}
    \item \textbf{Caminho (\textit{path}):} sequência de vértices $\langle v_0, v_1, \dots, v_k \rangle$ tal que $(v_{i-1}, v_i) \in E$ para todo $i \in \{1, \dots, k\}$;
    \item \textbf{Ciclo:} caminho fechado no qual $v_0 = v_k$ e que contém pelo menos uma aresta distinta. Um grafo que não possui ciclos é dito acíclico;
    \item \textbf{Grafo Acíclico Direcionado (DAG --- \textit{Directed Acyclic Graph}):} dígrafo desprovido de caminhos orientados fechados, estrutura fundamental para expressar relações de precedência estrita e causalidade \cite{cormen2009introduction, wang2026designerlyloop};
    \item \textbf{Conectividade e Conectividade Forte:} em grafos não direcionados, a conectividade assegura que exista ao menos um caminho entre qualquer par de vértices. Em dígrafos, define-se conectividade forte quando, para todo par de vértices $u, v \in V$, existem simultaneamente caminhos direcionados de $u$ para $v$ e de $v$ para $u$ \cite{cormen2009introduction};
    \item \textbf{Alcançabilidade (\textit{reachability}):} condição na qual existe um caminho orientado viável partindo de um nó inicial $v_{\text{start}}$ até um nó de destino $v_{\text{target}}$ \cite{cormen2009introduction, xu2025databasedrivenframework3dlevel};
    \item \textbf{Planaridade:} propriedade geométrica de um grafo que permite desenhá-lo em um plano bidimensional $\mathbb{R}^2$ de modo que suas arestas se interceptem unicamente nos vértices aos quais estão conectadas \cite{cormen2009introduction, linden2013designing}.
\end{itemize}

\subsubsection{Grafos Aplicados ao \textit{Level Design}}
\label{sec:grafos_level_design}

Na modelagem de fases e masmorras procedurais, a literatura apoia-se em grafos sob duas perspectivas complementares: grafos de espaço (topológicos) e grafos de missão (ou de ações) \cite{dormans2010, linden2013designing, linden2014dungeons, summerville2018procedural}.

Nos \textit{grafos topológicos de espaço}, cada vértice $v \in V$ representa uma região física discreta do nível (como uma sala, corredor, câmara ou pátio), e cada aresta $e \in E$ representa uma passagem física ou conexão de adjacência (portas, arcos de passagem, escadarias ou pontes) \cite{dormans2010, linden2014dungeons, summerville2018procedural}. Os vértices recebem atributos temáticos e funcionais (sala de início, sala de combate, câmara de tesouro, sala de chefe), enquanto as arestas recebem atributos de passagem (passagem aberta, porta trancada, passagem unidirecional com queda) \cite{linden2014dungeons, xu2025databasedrivenframework3dlevel}.

Nos \textit{grafos de missão} ou \textit{grafos de ações} (\textit{action graphs}), por outro lado, os vértices representam eventos, tarefas ou ações que o jogador executa (ex.: obter chave, ativar alavanca, derrotar sentinela, transpor portal), e as arestas direcionadas codificam restrições causais de precedência no tempo \cite{dormans2010, linden2013designing}. Conforme demonstrado seminalmente por Dormans \cite{dormans2010} e aprofundado por van der Linden et al. \cite{linden2013designing, linden2014dungeons}, um grafo de ações pode ser mapeado em um grafo de salas mediante o agrupamento de ações que compartilham a mesma localização física (\textit{co-located actions}), unificando a estrutura de progressão e a topologia arquitetônica.

\subsubsection{Grafos de Progressão e Dependências}
\label{sec:grafos_progressao}

A integração formal entre a conectividade espacial e as mecânicas de jogo estabelece os \textit{grafos de progressão}, nos quais arestas e vértices são condicionados a regras de acesso lógico \cite{dormans2010, linden2013designing, linden2014dungeons, xu2025databasedrivenframework3dlevel}. O paradigma central para estruturar a progressão em masmorras e fases de aventura é o quebra-cabeça de chave e fechadura (\textit{lock-and-key puzzle}) \cite{dormans2010, linden2014dungeons, xu2025databasedrivenframework3dlevel}. Uma fechadura denota qualquer barreira ou pré-requisito funcional (uma porta física trancada, uma barreira mágica, um obstáculo ambiental intransponível), enquanto a chave corresponde ao elemento ou estado necessário para desativá-la (uma chave física, uma habilidade de salto adquirida, o acionamento de um botão remoto ou a eliminação de inimigos) \cite{dormans2010, linden2013designing, shaker2016constructive}.

Do ponto de vista formal, a estrutura de dependências lógicas induz um subgrafo direcionado de pré-condições $G_{\text{dep}} = (V_{\text{dep}}, E_{\text{dep}})$, o qual deve satisfazer estritamente a condição de ser um Grafo Acíclico Direcionado (DAG) para garantir a inexistência de impasses circulares insolúveis (\textit{deadlocks}) \cite{cormen2009introduction, xu2025databasedrivenframework3dlevel}. Em tal DAG de dependências, existe ao menos uma ordenação topológica $\tau: V_{\text{dep}} \to \{1, 2, \dots, |V_{\text{dep}}|\}$ tal que, para qualquer aresta orientada $(u, v) \in E_{\text{dep}}$, tem-se $\tau(u) < \tau(v)$ \cite{cormen2009introduction}. Isso significa que todo recurso ou chave $u$ indispensável para a superação do bloqueio em $v$ deve ser alcançável e coletado antes que o estado de travessia de $v$ seja solicitado \cite{xu2025databasedrivenframework3dlevel}.

É essencial pontuar o contraste estrutural entre a lógica de dependências e a circulação física: enquanto o grafo causal de pré-condições ($G_{\text{dep}}$) deve ser estritamente acíclico para ser solucionável, o grafo espacial de navegação do nível pode conter ciclos controlados \cite{dormans2010, linden2014dungeons}. Esses ciclos na geometria física enriquecem a exploração ao fornecer rotas alternativas, atalhos destraváveis e a necessidade intencional de retorno a áreas previamente visitadas (\textit{backtracking}) após a obtenção de novos recursos \cite{dormans2010, linden2013designing, linden2014dungeons}. A solvabilidade do nível modela-se formalmente no espaço de estados do jogo ($S$), no qual cada estado $s \in S$ é composto pela localização espacial do agente associada ao seu inventário corrente de chaves, itens e modificadores de mundo \cite{dormans2010, linden2014dungeons, xu2025databasedrivenframework3dlevel}.

\subsubsection{Estrutura Lógica versus Representação Espacial}
\label{sec:logica_vs_espacial}

A Teoria dos Grafos elucida a separação ontológica entre a estrutura lógica e a representação geométrica: um grafo de progressão é uma entidade abstrata, topológica e combinatória, e não uma especificação espacial métrica \cite{dormans2010, summerville2018procedural}. Seus vértices e arestas estipulam a conectividade e a causalidade entre elementos, mas não contêm informações a respeito de posições tridimensionais $(x, y, z)$, comprimentos em metros, geometrias de paredes ou desníveis verticais \cite{dormans2010, linden2013designing, linden2014dungeons}.

Essa dissociação estabelece uma relação de mapeamento de um-para-muitos: um mesmo grafo de progressão admite infinitas configurações espaciais distintas no espaço contínuo tridimensional, contanto que as adjacências e ordenações lógicas sejam satisfeitas \cite{dormans2010, linden2013designing}. Todavia, nem todo grafo lógico abstrato é imediatamente realizável em uma grade ortogonal ou em um plano arquitetônico contínuo \cite{linden2013designing, linden2014dungeons}. Se um vértice do grafo contiver um número de arestas incidentes (grau) superior ao número físico de portas comportadas pela geometria do módulo arquitetônico, ou se a topologia for não planar no plano bidimensional, a instanciação física necessita de técnicas de planarização e expansão estrutural, tais como a introdução algorítmica de nós intermediários (corredores ramificados, bifurcações ou transições verticais entre múltiplos pavimentos), conforme formalizado por van der Linden et al. \cite{linden2013designing, linden2014dungeons}.

% ---------------------------------------------------------------------------
\subsection{Representação de \textit{Layouts} Espaciais 3D}
\label{sec:layout_3d}
% ---------------------------------------------------------------------------

\subsubsection{Representação Modular de Ambientes}
\label{sec:representacao_modular}

Para converter a especificação abstrata de um grafo em um ambiente interativo tridimensional, a literatura de PCG e os motores de jogos apoiam-se no paradigma da modelagem espacial modular \cite{linden2014dungeons, shaker2016constructive, xu2025databasedrivenframework3dlevel}. Nessa abordagem, o espaço não é concebido como uma malha estática singular, mas através da composição combinatória de unidades tridimensionais reutilizáveis e pré-autoradas \cite{linden2014dungeons, shaker2016constructive, xu2025databasedrivenframework3dlevel}.

As granularidades modulares típicas abrangem \cite{linden2014dungeons, shaker2016constructive, vieira2024metodologia}:
\begin{itemize}
    \item \textbf{Grades regulares e \textit{voxels}:} partição do espaço em células discretas uniformes em 2D ou 3D, atribuindo-se a cada célula um elemento primitivo de geometria (bloco sólido, piso, parede, abertura ou ar livre) \cite{pcgbook, shaker2016constructive, summerville2018procedural};
    \item \textbf{Segmentos compostos (\textit{chunks}):} subestruturas volumétricas que agrupam conjuntos pré-definidos de blocos e acessórios temáticos \cite{linden2014dungeons, shaker2016constructive};
    \item \textbf{Módulos de salas e \textit{prefabs} 3D:} blocos arquitetônicos inteiros (como salas com dimensões estipuladas, corredores em L, junções em T e escadarias) e ativos pré-fabricados (\textit{prefabs}) de objetos internos \cite{linden2014dungeons, xu2025databasedrivenframework3dlevel, regmi2025system}. Um \textit{prefab} encapsula a malha gráfica 3D (\textit{mesh}), materiais, texturas, colunas de suporte, volumes de colisão física e pontos de ancoragem semânticos em um único pacote modular \cite{xu2025databasedrivenframework3dlevel, regmi2025system}.
\end{itemize}

O acoplamento topológico e geométrico entre módulos adjacentes é operacionalizado por meio de conectores ou portas (\textit{connectors} / \textit{doorways} / \textit{sockets}) posicionados estrategicamente nas extremidades dos blocos modulares \cite{linden2013designing, linden2014dungeons, shaker2016constructive, sun2024roomplanner}. Cada conector define uma posição métrica relativa em relação ao pivô do módulo, a dimensão da abertura e um vetor de orientação normal que especifica a direção para a qual a passagem se projeta no espaço \cite{linden2013designing, shaker2016constructive, sun2024roomplanner, xu2025databasedrivenframework3dlevel}.

\subsubsection{Relações e Transformações Espaciais}
\label{sec:transformacoes_espaciais}

A instanciação de peças modulares no espaço contínuo $\mathbb{R}^3$ fundamenta-se nas transformações afins da álgebra linear, articulando o referencial de coordenadas local de cada objeto e o referencial global do cenário (\textit{world space}) \cite{ericson2004realtime, sun2024roomplanner, xu2025constraint}.

A especificação espacial completa de um componente $i$ exige distinguir sua \textit{pose}, suas dimensões físicas e seus fatores de escala \cite{ericson2004realtime, sun2024roomplanner, xu2025constraint}:
\begin{itemize}
    \item A \textbf{pose} define a localização e a atitude espacial do componente no mundo, decomposta na translação $\mathbf{p}_i = (x_i, y_i, z_i) \in \mathbb{R}^3$ (posição de seu ponto de pivô) e na rotação $\mathbf{R}_i \in SO(3)$ (orientação angular, expressa por matrizes ortogonais, quatérnios ou ângulos de Euler, tipicamente parametrizada em layouts pelo ângulo de guinada $\theta_i$ em torno do eixo vertical $Y$ para rotações ortogonais de $90^\circ$ ou $180^\circ$) \cite{sun2024roomplanner, xu2025databasedrivenframework3dlevel};
    \item As \textbf{dimensões nominais} $(w_i, l_i, h_i) \in \mathbb{R}^3$ representam a largura, o comprimento e a altura intrínsecos do volume do objeto no espaço local;
    \item O \textbf{fator de escala} $\mathbf{s}_i = (s_{x,i}, s_{y,i}, s_{z,i}) \in \mathbb{R}^3$ expressa a ampliação ou redução proporcional aplicada à geometria nos eixos coordenados.
\end{itemize}

Para a detecção de colisões e verificação de proximidade geométrica em tempo computacional eficiente, a computação gráfica recorre a volumes delimitadores (\textit{bounding volumes}) \cite{ericson2004realtime, sun2024roomplanner, xu2025constraint}. Distinguem-se fundamentalmente dois tipos:
\begin{enumerate}
    \item \textbf{Caixa Delimitadora Alinhada aos Eixos (\textit{Axis-Aligned Bounding Box} --- AABB):} definida pelos intervalos extremos $[x_{\min}, x_{\max}] \times [y_{\min}, y_{\max}] \times [z_{\min}, z_{\max}]$ paralelamente aos eixos globais do sistema de coordenadas \cite{ericson2004realtime}. Testes de interseção entre AABBs possuem custo computacional extremamente baixo por meio de comparações escalares de intervalos em cada eixo; contudo, quando um objeto sofre rotação angular arbitrária, sua AABB precisa ser recalculada e expandida para envolver os novos limites, introduzindo espaço vazio excessivo que pode gerar falsos positivos de colisão \cite{ericson2004realtime};
    \item \textbf{Caixa Delimitadora Orientada (\textit{Oriented Bounding Box} --- OBB):} paralelepípedo delimitador centrado no objeto cujos eixos permanecem solidários à rotação e orientação da própria geometria local \cite{ericson2004realtime}. Embora exija testes de separação mais sofisticados (como o Teorema dos Eixos Separadores --- SAT), a OBB ajusta-se com precisão estrita ao volume sólido sob qualquer rotação espacial \cite{ericson2004realtime}.
\end{enumerate}

\subsubsection{Restrições Geométricas}
\label{sec:restricoes_geometricas}

A montagem de um \textit{layout} tridimensional coerente requer a satisfação concorrente de um conjunto estrito de restrições geométricas \cite{linden2014dungeons, ericson2004realtime, sun2024roomplanner, xu2025constraint, xu2025databasedrivenframework3dlevel}:
\begin{itemize}
    \item \textbf{Restrição de não sobreposição volumétrica (\textit{non-overlap}):} sólidos arquitetônicos e objetos distintos não podem ocupar concomitantemente a mesma porção do espaço euclidiano tridimensional. Formalmente, a interseção volumétrica entre seus interiores sólidos deve ser estritamente nula ($V_{\text{overlap}} = 0$ ou $\text{IoU}_{3D} = 0$), evitando interpenetrações anômalas (\textit{clipping}) \cite{ericson2004realtime, sun2024roomplanner, xu2025constraint};
    \item \textbf{Restrição de contenção e fronteiras (\textit{boundary constraints}):} objetos internos associados a uma sala devem ter seus volumes contidos estritamente no interior do polígono ou prisma formado por suas paredes, piso e teto, sem atravessar a alvenaria perimetral \cite{sun2024roomplanner, xu2025constraint};
    \item \textbf{Compatibilidade e alinhamento de portais:} para que duas salas ou uma sala e um corredor compartilhem uma passagem transitável, seus respectivos conectores devem coincidir métrica e volumetricamente na fronteira comum, apresentando vetores normais opostos e colineares, sem que ocorra bloqueio por paredes cegas de um dos módulos \cite{linden2013designing, shaker2016constructive, sun2024roomplanner, xu2025databasedrivenframework3dlevel};
    \item \textbf{Desobstrução e espaço de circulação (\textit{clearance}):} a distribuição de objetos em um cômodo deve resguardar zonas livres desimpedidas ao redor de portas, escadas e rotas centrais de passagem, prevenindo o congestionamento ou adensamento físico indevido (\textit{overcrowding} / \textit{clutter}) que impeça o trânsito físico dos agentes virtuais \cite{sun2024roomplanner, xu2025constraint, xu2025databasedrivenframework3dlevel}.
\end{itemize}

% ---------------------------------------------------------------------------
\subsection{Arquiteturas Híbridas e Geração Baseada em Restrições}
\label{sec:arquiteturas_hibridas}
% ---------------------------------------------------------------------------

\subsubsection{Sistemas Híbridos e Neuro-Simbólicos}
\label{sec:sistemas_hibridos}

A análise aprofundada dos paradigmas generativos revela uma complementaridade ontológica: métodos puramente simbólicos, construtivos ou orientados a restrições fornecem rigor matemático e preservação determinística de invariantes funcionais, porém demandam complexa formalização manual de regras e tendem a apresentar menor flexibilidade estilística \cite{shaker2016constructive, summerville2018procedural}; inversamente, modelos baseados em redes profundas e LLMs exibem expressividade de linguagem, compreensão semântica e facilidade na geração de conteúdo criativo, porém sofrem de alucinações espaciais e ausência de garantias formais de consistência \cite{Gallotta_2024, maleki2024procedural}.

Essa constatação catalisou a emergência de \textit{arquiteturas híbridas neuro-simbólicas} em inteligência artificial e geração de fases \cite{garcez2023neurosymbolic, Gallotta_2024, xu2025constraint}. Conforme teorizado por Garcez e Lamb \cite{garcez2023neurosymbolic}, o paradigma neuro-simbólico postula a integração sinérgica entre o aprendizado conexionista baseado em dados (capacidade de inferência intuitiva, correlação estatística e semântica de alto nível) e o raciocínio simbólico (estruturas lógicas, grafos, solucionadores e verificação formal baseada em regras). No contexto de PCG para ambientes 3D, essa divisão funcional traduz-se na delegação das etapas de planejamento conceitual, estruturação lógica e contextualização ao módulo de linguagem (neural), ao passo que as rotinas de resolução métrica contínua, instanciação modular e validação geométrica são executadas por algoritmos determinísticos e solucionadores formais (simbólicos) \cite{hu2024scenecraft, tang2025unrealllm, chen2025narrative, xu2025constraint, xu2025databasedrivenframework3dlevel, her2025zeroshot3dmapgeneration}.

De modo análogo, essa divisão conceitual dialoga com os princípios da geração de iniciativa mista (\textit{Mixed-Initiative PCG}), fundamentada seminalmente em ferramentas como \textit{Tanagra} \cite{tanagra} e \textit{Sentient Sketchbook} \cite{liapis2013sentient}, nas quais o sistema computacional assume o papel de parceiro proativo que calcula métricas espaciais em tempo de autoria, resolve restrições e sugere correções estruturais mantendo o controle deliberativo do designer \cite{tanagra, liapis2013sentient, lai2020towards, walton2020evaluating, wang2026designerlyloop}.

\subsubsection{Geração Baseada em Regras e Restrições}
\label{sec:geracao_restricoes}

O suporte matemático para a imposição de regras em sistemas generativos reside na teoria de Problemas de Satisfação de Restrições (CSP --- \textit{Constraint Satisfaction Problem}) \cite{russell2021artificial}. Formalmente, um CSP é definido pela tupla $(X, D, C)$, na qual $X = \{X_1, X_2, \dots, X_n\}$ denota um conjunto de variáveis de decisão (como coordenadas espaciais, orientações ou tipos de sala), $D = \{D_1, D_2, \dots, D_n\}$ representa os domínios de valores admissíveis para cada variável, e $C = \{C_1, C_2, \dots, C_m\}$ constitui um conjunto de restrições que delimitam as combinações simultâneas de valores permitidas \cite{russell2021artificial}. A resolução de um CSP emprega algoritmos de retrocesso sistemático (\textit{backtracking}) potencializados por heurísticas de ordenação e mecanismos de propagação de restrições, a exemplo da consistência de arcos (\textit{arc consistency}), que elimina antecipadamente valores de domínios que comprovadamente violariam condições futuras \cite{russell2021artificial}.

No escopo do arranjo espacial e planejamento de fases, as restrições são categorizadas em duas ordens de rigor \cite{tanagra, xu2025constraint, xu2025databasedrivenframework3dlevel}:
\begin{itemize}
    \item \textbf{Restrições rígidas (\textit{hard constraints}):} condições mandatórias e invioláveis para a viabilidade do nível, cuja inobservância inviabiliza funcional ou fisicamente o artefato (tais como ausência de sobreposição sólida, contenção no perímetro e ordenação causal acíclica de chaves e portas) \cite{tanagra, sun2024roomplanner, xu2025constraint};
    \item \textbf{Restrições flexíveis (\textit{soft constraints}):} diretrizes de qualidade, preferências estéticas ou métricas de balanceamento que expressam objetivos desejáveis (tais como distâncias médias de deslocamento, agrupamento estilístico de móveis ou dispersão de desafios), cujo atendimento é otimizado sem configurar critério binário de rejeição \cite{xu2025constraint, xu2025databasedrivenframework3dlevel}.
\end{itemize}

Quando o problema de disposição espacial envolve a ponderação conjunta de restrições rígidas e flexíveis em domínios contínuos, ele é frequentemente modelado por analogia ao clássico Problema de \textit{Layout} de Instalações (\textit{Facility Layout Problem} --- FLP) da pesquisa operacional \cite{xu2025constraint, xu2025databasedrivenframework3dlevel}. Sob essa formulação, define-se uma função objetivo com penalidades ponderadas associadas a cada infração de distância ou posicionamento, buscando configurações que minimizem o custo global de violações \cite{xu2025constraint}.

\subsubsection{Validação e Reparo de Conteúdo Gerado}
\label{sec:validacao_reparo}

Dado que geradores conexionistas, estocásticos ou heurísticos não garantem a priori a total conformidade com restrições rígidas, a literatura formaliza quatro paradigmas teóricos para a garantia da correção do conteúdo \cite{pcgbook, summerville2018procedural, maleki2024procedural, xu2025databasedrivenframework3dlevel}:
\begin{enumerate}
    \item \textbf{Geração e teste com filtragem (Rejeição integral):} o gerador sintetiza o candidato e o submete a uma bateria de verificações algorítmicas; constatada qualquer infração rígida, o candidato é sumariamente descartado e o procedimento reiniciado \cite{pcgbook, maleki2024procedural};
    \item \textbf{Construção com retrocesso (\textit{Backtracking}):} o sistema aloca componentes de forma incremental; identificada a impossibilidade de satisfação futura das restrições no nó atual da árvore de busca, o algoritmo desfaz a última atribuição e explora caminhos alternativos \cite{tanagra, linden2014dungeons, russell2021artificial};
    \item \textbf{Reparo pós-geração (\textit{Post-generation repair}):} operadores algorítmicos locais atuam sobre o artefato gerado para solucionar violações pontuais sem descartá-lo, tais como transladar pequenos móveis em colisão, regenerar conectores desalinhados ou desobstruir passagens bloqueadas \cite{summerville2018procedural, sun2024roomplanner, xu2025databasedrivenframework3dlevel};
    \item \textbf{Refinamento iterativo em malha fechada (\textit{Feedback loop}):} diagnósticos detalhados extraídos pelos validadores formais são formatados como mensagens estruturadas de erro e reinjetados no módulo generativo de alto nível, permitindo correções informadas em iterações subsequentes \cite{hu2024scenecraft, nasir2024word2worldgeneratingstoriesworlds, xu2025constraint, her2025zeroshot3dmapgeneration}.
\end{enumerate}

% ---------------------------------------------------------------------------
\subsection{Validade, Navegabilidade e Solvabilidade de Fases}
\label{sec:validade_navegabilidade_solvabilidade}
% ---------------------------------------------------------------------------

A avaliação formal da correção de uma fase de jogo exige decompor rigorosamente o conceito genérico de ``jogabilidade funcional'' em quatro dimensões teóricas distintas: validade lógica, validade geométrica, navegabilidade física e solvabilidade global. Essas dimensões operam em planos de abstração distintos e não configuram sinônimos \cite{dormans2010, linden2013designing, linden2014dungeons, sun2024roomplanner, xu2025constraint}.

\subsubsection{Validade Lógica e Estrutural}
\label{sec:validade_logica}

A validade lógica e estrutural diz respeito à consistência relacional e causal do nível em sua representação topológica abstrata, desconsiderando dimensões métricas ou físicas do espaço tridimensional \cite{dormans2010, linden2013designing, linden2014dungeons}. Uma fase é considerada logicamente válida quando satisfaz as invariantes fundamentais do sistema de regras:
\begin{itemize}
    \item Presença de nós de início e fim devidamente parametrizados;
    \item Conectividade do grafo subjacente, atestando que todas as salas e objetivos necessários integram a mesma componente topológica conexa \cite{linden2014dungeons, xu2025databasedrivenframework3dlevel};
    \item Consistência e aciclicidade das dependências causais ($G_{\text{dep}}$ é um DAG), assegurando que não existem bloqueios circulares insolúveis, tais como trancar o acesso à chave de uma porta atrás da própria porta trancada \cite{dormans2010, linden2014dungeons, xu2025databasedrivenframework3dlevel}.
\end{itemize}

\subsubsection{Validade Geométrica}
\label{sec:validade_geometrica}

A validade geométrica refere-se à integridade física do ambiente quando seus módulos são instanciados no espaço contínuo tridimensional $\mathbb{R}^3$ \cite{ericson2004realtime, sun2024roomplanner, xu2025constraint, feng2026repurposing3dgenerativemodel}. Um layout 3D possui validade geométrica se, e somente se:
\begin{itemize}
    \item Nenhuma malha arquitetônica sólida (paredes, pisos ou mobílias densas) sobrepõe-se a outra com volume de interseção positivo ($V_{\text{overlap}} = 0$) \cite{ericson2004realtime, xu2025constraint};
    \item Todos os elementos internos respeitam os limites envoltórios do compartimento que os abriga \cite{sun2024roomplanner, xu2025constraint};
    \item Todas as portas e conectores adjacentes previstos na topologia estão devidamente alinhados e desimpedidos por paredes estruturais \cite{linden2013designing, sun2024roomplanner, xu2025databasedrivenframework3dlevel}.
\end{itemize}

A verificação dessa dimensão em motores gráficos e sistemas computacionais apoia-se em estruturas espaciais de consulta (como AABBs, OBBs, octrees e algoritmos de varredura como \textit{raycasting}) \cite{shaker2016constructive, ericson2004realtime, sun2024roomplanner}. Sob a ótica teórica, essas estruturas operam como ferramentas determinísticas de detecção e consulta espacial, fornecendo subsídios para que algoritmos de verificação confirmem ou rejeitem a conformidade geométrica do layout \cite{shaker2016constructive, xu2025constraint}.

\subsubsection{Navegabilidade e Malhas de Navegação}
\label{sec:navegabilidade}

A navegabilidade (\textit{navigability} ou \textit{traversability}) atesta se o ambiente tridimensional oferece superfícies contínuas e transponíveis sobre as quais um agente cinemático (jogador ou NPC) é fisicamente capaz de se locomover \cite{snook2000simplified, sun2024roomplanner, xu2025constraint, xu2025databasedrivenframework3dlevel}. A navegabilidade é uma propriedade dinâmica e dependente dos atributos físicos do agente, tais como altura livre requerida, raio cilíndrico de colisão, inclinação máxima de superfícies e altura máxima de degraus superáveis \cite{snook2000simplified, xu2025constraint}. Um nível pode exibir perfeita validade geométrica (ausência de colisão entre malhas), mas ser desprovido de navegabilidade caso um corredor seja mais estreito que o diâmetro do agente ou caso ocorra um desnível vertical desprovido de escada ou rampa de acesso \cite{sun2024roomplanner, xu2025databasedrivenframework3dlevel}.

A representação computacional padrão para o espaço contínuo transitável em ambientes tridimensionais consiste na Malha de Navegação (\textit{Navigation Mesh} --- NavMesh), introduzida na literatura de jogos por Snook \cite{snook2000simplified}. A NavMesh discretiza as superfícies caminháveis (\textit{walkable areas}) da geometria 3D decompondo-as em um conjunto de polígonos convexos bidimensionais interconectados no espaço tridimensional \cite{snook2000simplified, regmi2025system}. A convexidade de cada polígono assegura que qualquer trajetória linear traçada entre dois pontos pertencentes ao mesmo polígono permanece estritamente contida dentro da área caminhável, sem colidir com obstáculos \cite{snook2000simplified}. O cálculo de rotas em uma NavMesh opera sobre o \textit{grafo dual de adjacência} formado pelos polígonos convexos (em que os polígonos constituem vértices e as arestas compartilhadas constituem arestas do grafo dual), aplicando-se posteriormente algoritmos de suavização geométrica (como o \textit{funnel algorithm} ou \textit{string pulling}) para gerar caminhos esteticamente suaves e métricamente otimizados \cite{snook2000simplified, xu2025databasedrivenframework3dlevel}.

\subsubsection{Busca de Caminhos (\textit{Pathfinding})}
\label{sec:pathfinding}

A busca de caminhos (\textit{pathfinding}) consiste no procedimento algorítmico responsável por determinar uma rota contínua de travessia entre uma posição inicial $\mathbf{p}_{\text{start}}$ e uma posição alvo $\mathbf{p}_{\text{target}}$ sobre o espaço transitável discretizado (como a NavMesh) \cite{hart1968formal, snook2000simplified, sun2024roomplanner, xu2025databasedrivenframework3dlevel}.

O algoritmo $A^*$ (\textit{A-star}), formulado por Hart, Nilsson e Raphael \cite{hart1968formal}, é a solução canônica para busca informada de caminhos em grafos de navegação. O $A^*$ avalia vértices candidatos por meio de uma função de custo heurístico:
\begin{equation}
    f(n) = g(n) + h(n)
\end{equation}
onde $g(n)$ corresponde ao custo exato do caminho já percorrido da origem até o nó $n$, e $h(n)$ denota a estimativa heurística de custo de $n$ até o objetivo \cite{hart1968formal}. Uma propriedade fundamental demonstrada por Hart et al. \cite{hart1968formal} reside no critério de admissibilidade: se a função heurística $h(n)$ for admissível --- ou seja, se nunca superestimar a distância real mínima até o alvo ($h(n) \le h^*(n)$, como ocorre com a distância euclidiana em linha reta no espaço contínuo) ---, o algoritmo $A^*$ tem garantia matemática de encontrar um caminho de custo ótimo (mínimo), caso exista ao menos uma rota transponível \cite{hart1968formal}.

No contexto da verificação algorítmica de fases, a execução de testes de $A^*$ entre pontos críticos comprova que o espaço modelado mantém conectividade física e operacional real entre compartimentos distintos \cite{sun2024roomplanner, xu2025databasedrivenframework3dlevel}.

\subsubsection{Solvabilidade de uma Fase}
\label{sec:solvabilidade}

A solvabilidade (\textit{solvability}, \textit{winnability} ou completabilidade funcional) constitui o critério avaliativo superior de uma fase: ela estipula a existência de pelo menos uma trajetória válida e realizável pelo jogador no espaço de estados e no espaço geométrico tridimensional, partindo da configuração inicial e atingindo a condição de vitória sem violar nenhuma restrição física ou mecânica do jogo \cite{tanagra, dormans2010, linden2014dungeons, summerville2018procedural, nasir2024word2worldgeneratingstoriesworlds, xu2025constraint}.

A distinção teórica entre validade geométrica, navegabilidade e solvabilidade é indispensável porque nenhuma dessas condições é suficiente para assegurar as demais \cite{dormans2010, linden2014dungeons, summerville2018procedural, xu2025constraint}:
\begin{enumerate}
    \item \textbf{Validade geométrica sem navegabilidade:} dois recintos podem ser posicionados sem nenhuma colisão volumétrica entre suas paredes, mas separados por uma vala intransponível ou por uma rampa com inclinação superior ao limite de locomoção do agente, inviabilizando a transposição física \cite{sun2024roomplanner, xu2025databasedrivenframework3dlevel};
    \item \textbf{Navegabilidade física sem solvabilidade dinâmica:} o algoritmo $A^*$ pode computar com sucesso um caminho contínuo na NavMesh estática entre a entrada e a saída do nível; entretanto, se esse trajeto interceptar uma porta trancada cuja respectiva chave estiver situada em uma câmara cujo acesso físico dependa da prévia transposição dessa mesma porta, o nível apresenta navegabilidade métrica, porém é absolutamente insolúvel devido ao impasse lógico de progressão \cite{dormans2010, linden2013designing, xu2025constraint};
    \item \textbf{Validade lógica abstrata sem solvabilidade espacial:} um grafo de progressão pode ter todas as suas relações de precedência perfeitamente ordenadas no plano abstrato ($\tau(chave) < \tau(porta)$), mas o nível resultante ser insolúvel porque a instanciação geométrica 3D posicionou um obstáculo físico denso que bloqueou o corredor de acesso à sala da chave, rompendo a transitabilidade real da fase \cite{linden2014dungeons, xu2025databasedrivenframework3dlevel}.
\end{enumerate}

Portanto, a solvabilidade plena em ambientes tridimensionais com progressão não linear requer que a alcançabilidade seja verificada simultaneamente sob duas dimensões acopladas: a validade causal no espaço de estados lógicos e a conectividade física contínua comprovada no espaço métrico navegável da fase \cite{dormans2010, linden2014dungeons, xu2025constraint, xu2025databasedrivenframework3dlevel}.

\subsubsection{Síntese Conceitual e Perspectiva Metodológica}
\label{sec:sintese_transicao}

A articulação dos conceitos expostos ao longo deste capítulo evidencia que a síntese de fases tridimensionais complexas impõe um duplo desafio formal: requer, por um lado, alta capacidade semântica e contextual para conceber narrativas, estruturas de desafio e progressões lógicas que regulem o ritmo do jogo; e exige, por outro lado, rigor determinístico absoluto para satisfazer restrições geométricas volumétricas, assegurar alinhamentos modulares e garantir a transitabilidade da malha de navegação.

O reconhecimento das limitações intrínsecas dos modelos conexionistas na manipulação direta de coordenadas métricas, contraposto ao poder estruturante dos formalismos de grafos e da satisfação de restrições, consolida a fundamentação teórica que orienta as abordagens modernas de cocriação e automação em jogos. Esse arcabouço conceitual --- alicerçado na estrita distinção entre organização lógica e realização espacial, e na decomposição analítica entre validade geométrica, navegabilidade e solvabilidade --- estabelece as premissas indispensáveis para a formulação dos modelos, representações intermediárias e rotinas determinísticas de instanciação que serão detalhados no capítulo de Metodologia.
